\section{Terminological provenance and bounded prior-use search}
The exploratory notebooks used the working expression ``useownership''. The present revision adopts \emph{usownership} as the noun for the generalised institution, but the associated \emph{Usowner} concept is not being backdated to 2026. ERC-4950, created on 28 March 2022 by V\'ictor Mu\~noz, Josep Lluis de la Rosa and Easy Innova, explicitly lists ``Usowners'' among its applications and glosses them as users who consume an NFT and become partial owners \citep{erc4950}. The same proposal implements a pair of entangled ERC-721 tokens that can jointly operate a smart-contract wallet. Its current status is \emph{Stagnant}; the manuscript therefore calls it a standards-track ERC proposal or technical precursor, not an adopted Ethereum standard.

The present noun is constructed from the same semantic transition: \emph{user} $\rightarrow$ \emph{Usowner} $\rightarrow$ \emph{usownership}, where \emph{-ship} denotes the durable ownership relation. What is new in the present paper is the generalisation from an NFT-oriented partial-owner use case to a recurring corporate mechanism with graded contribution, dilution accounting and macroeconomic transition dynamics. The 2022 ERC does not contain the contribution-score equation, automation-linked issuance schedule, matched fiscal benchmark or monetary-financing model used here. This distinction is important because it lets the paper claim continuity without claiming novelty for its own precursor.

The prior-use check was designed to avoid an unsupported lexical claim rather than to establish trademark clearance. On 9 September 2026 exact-string searches were run for ``usownership'' and variants combined with \emph{ownership}, \emph{user ownership}, \emph{platform}, \emph{cooperative}, \emph{patent}, \emph{trademark}, \emph{book} and \emph{dictionary}; additional queries were scoped to academic-web surfaces. The search found unrelated labels such as \texttt{usOwnership} meaning U.S. ownership \citep{straitsx2026}, but no indexed source using \emph{usownership} as the noun for the generalised contribution-weighted user-to-owner institution defined here. The claim remains semantic and mechanism-specific, can be overturned by unindexed prior use, and is not a legal priority opinion.

Conceptual priority is also explicitly bounded. Dunn's 1988 cooperative principles place user, owner and control roles in the same individual \citep{dunn1988}; platform cooperativism later described user-owned and produser-owned platforms \citep{scholz2016}; and \citet{schneider2018} analysed an Internet of ownership. The distinctive object tested here is gradual, recurring, contribution-weighted equity acquisition inside an initially investor-owned firm, linked to automation intensity, income coverage and monetary constraints.

\section{Prior Internet-of-Value research lineage}
De la Rosa Esteva's chapter \emph{Potential Sources of Internet of Value Systemic Risk} appeared in the Springer volume \emph{Enabling the Internet of Value} in 2022, pp.~185--188 \citep{delarosa2022iov}. The chapter argues that lack of interoperability and scalability can become systemic risks as tokenised claims traverse networks, links IoV business models to the democratisation of finance and property, and anticipates a ``massive migration of value onto the ledgers'' described as a \emph{Value Deluge}. It further argues that curation, integrity and ownership management are necessary to preserve the value of digital assets across exchanges. This is the direct architectural precursor used in the main text.

The exact chapter should not be cited for claims it does not make. It does not formalise the current heterogeneous contribution score or derive ownership shares from measured user activity. Other chapters in the wider IoV volume discuss machine-readable usage, usage licences, telemetry, royalties and micro-licensing, but those contributions have different authors. In the present research lineage, ERC-4950 is the source that explicitly states the partial-user-owner idea, while this paper supplies the general contribution-measurement and corporate-equity formalisation. Keeping those roles separate makes the novelty account stronger and auditable.

\section{Internet-of-Value implementation boundary}
The expression ``Internet of Value'' predates both sources above. The academic literature has since developed broader definitions around tokenisation, digital platforms and smart contracts \citep{lohmer2024}. In 2026 the European Commission described DLT and tokenisation as potentially paving the way to an Internet of Value \citep{eciov2026}. BIS proposals for tokenisation and unified ledgers are particularly relevant because they embed programmable assets in trusted monetary architecture rather than requiring a permissionless cryptocurrency system \citep{bis2023,bis2025,bis2026}. The ECB reports that European DLT-based fixed-income issuance has reached close to EUR~4 billion since 2021 and that Eurosystem trials and experiments involved roughly EUR~1.6 billion across nine jurisdictions \citep{cipollone2026}. These figures establish implementation momentum, not proof that tokenisation will dominate financial markets or that it lowers the total social cost of usownership.

The main paper therefore uses the IoV as an emerging implementation direction. If digital ownership claims and payment settle on interoperable programmable rails, the marginal administrative cost of recurring participant equity may fall. No numerical cost saving is imposed in the simulations; Equation~\ref{M-eq:implementation} in the main manuscript states the break-even condition. Legal validity, contribution truth, privacy, cyber resilience and macroeconomic incidence remain outside what a ledger can solve by itself.
